\documentclass[10pt,a4paper]{article}
\usepackage[utf8]{inputenc}
\usepackage[vietnamese]{babel}
\usepackage{amsmath,amssymb,amsfonts}
\usepackage{geometry}
\geometry{a4paper, left=10mm, right=10mm, top=10mm, bottom=10mm}
\usepackage{tikz}
\usetikzlibrary{shapes,arrows.meta,positioning,calc}
\usepackage[most]{tcolorbox}
\usepackage{array,booktabs,multicol}
\usepackage{fontawesome}
\usepackage{helvet}
\renewcommand{\familydefault}{\sfdefault}
\usepackage{parskip}

\definecolor{primary}{RGB}{112, 48, 160}
\definecolor{teal}{RGB}{0, 128, 128}
\definecolor{orange}{RGB}{230, 108, 44}
\definecolor{green}{RGB}{46, 139, 87}
\definecolor{rose}{RGB}{192, 0, 0}

\pagestyle{empty}

\begin{document}

% --- HEADER BANNER ---
\begin{tcolorbox}[
  enhanced,
  colback=violet!8!white,
  colframe=violet!80!black,
  arc=2.5mm, boxrule=1pt,
  left=3mm, right=3mm, top=2mm, bottom=2mm
]
  \begin{tabular*}{\textwidth}{@{\extracolsep{\fill}} lr}
    {\small\bfseries\color{violet!85!black}\faBook\ SỔ TAY KIẾN THỨC TOÁN 2} & 
    {\small\bfseries\color{teal!85!black}\faGraduationCap\ CHƯƠNG TRÌNH GDPT 2018}
  \end{tabular*}
  \vspace{-1mm}
  \begin{center}
    \rule{0.98\textwidth}{0.6pt}
  \end{center}
  \vspace{-2mm}
  \centering
  {\LARGE\bfseries\color{violet!90!black}\faBookmark\ BÀI 59: PHÉP CỘNG (KHÔNG NHỚ) TRONG PHẠM VI 1000}\\[1mm]
  {\itshape\footnotesize\color{black!75} Kĩ thuật đặt tính và tính cộng các số có ba chữ số không nhớ}
\end{tcolorbox}

\vspace{-2mm}

\noindent
\begin{minipage}[t]{0.49\textwidth}\vspace{0pt}
\begin{tcolorbox}[
  enhanced,
  colback=violet!4!white,
  colframe=violet!75!black,
  arc=2mm, boxrule=0.8pt,
  left=2mm, right=2mm, top=1.8mm, bottom=1.8mm,
  title={\footnotesize\bfseries\color{white}\faKey\ 1. Quy tắc đặt tính và tính}
]
{\footnotesize
\begin{itemize}\setlength{\itemsep}{1pt}\setlength{\leftskip}{-2.5mm}
  \item \textbf{Bước 1: Đặt tính:}\\
  Viết số hạng này dưới số hạng kia sao cho:
  \begin{itemize}\setlength{\itemsep}{0.5pt}\setlength{\leftskip}{-2mm}
    \item Đơn vị thẳng cột với đơn vị.
    \item Chục thẳng cột với chục.
    \item Trăm thẳng cột với trăm.
    \item Viết dấu cộng $+$ bên trái, kẻ vạch ngang dưới hai số.
  \end{itemize}
  \item \textbf{Bước 2: Tính:}\\
  Thực hiện cộng lần lượt từ \textbf{phải sang trái}:
  \begin{itemize}\setlength{\itemsep}{0.5pt}\setlength{\leftskip}{-2mm}
    \item Cộng các chữ số hàng đơn vị.
    \item Cộng các chữ số hàng chục.
    \item Cộng các chữ số hàng trăm.
  \end{itemize}
\end{itemize}
}
\end{tcolorbox}
\end{minipage}
\hfill
\begin{minipage}[t]{0.49\textwidth}\vspace{0pt}
\begin{tcolorbox}[
  enhanced,
  colback=teal!4!white,
  colframe=teal!75!black,
  arc=2mm, boxrule=0.8pt,
  left=2mm, right=2mm, top=1.8mm, bottom=1.8mm,
  title={\footnotesize\bfseries\color{white}\faSliders\ 2. Phân tích chi tiết ví dụ mẫu}
]
{\footnotesize
\textbf{Tính: $352 + 416 = ?$}\\
\begin{itemize}\setlength{\itemsep}{1pt}\setlength{\leftskip}{-2.5mm}
  \item \textbf{Đặt tính:}
  \begin{center}
  \begin{tabular}{r@{\,}c@{\,}l}
    & $352$ & \\
  $+$ & $416$ & \\
  \hline
    & $\mathbf{768}$ &
  \end{tabular}
  \end{center}
  \item \textbf{Thực hiện tính:}
  \begin{itemize}\setlength{\itemsep}{0.5pt}\setlength{\leftskip}{-2mm}
    \item $2 + 6 = \mathbf{8}$, viết $8$ ở hàng đơn vị.
    \item $5 + 1 = \mathbf{6}$, viết $6$ ở hàng chục.
    \item $3 + 4 = \mathbf{7}$, viết $7$ ở hàng trăm.
  \end{itemize}
  \item Vậy: $\mathbf{352 + 416 = 768}$.
\end{itemize}
}
\end{tcolorbox}
\end{minipage}

\vspace{-1.5mm}
\begin{tcolorbox}[
  enhanced,
  colback=orange!4!white,
  colframe=orange!80!black,
  arc=2mm, boxrule=0.8pt,
  left=2mm, right=2mm, top=1.5mm, bottom=1.5mm,
  title={\footnotesize\bfseries\color{white}\faSitemap\ 3. Sơ đồ thứ tự thực hiện phép tính từ phải sang trái}
]
\centering
\begin{tikzpicture}[x=1.65cm, y=0.75cm]
  \tikzset{arr/.style={-{Stealth[scale=0.8]}}, every node/.append style={align=center}}
  \node[rectangle, rounded corners=2pt, draw=green!70!black, fill=green!15, font=\scriptsize\bfseries] (dv) at (0,0) {1. HÀNG ĐƠN VỊ\\$2 + 6 = 8$};
  \node[rectangle, rounded corners=2pt, draw=teal!80!black, fill=teal!15, font=\scriptsize\bfseries] (ch) at (2.4,0) {2. HÀNG CHỤC\\$5 + 1 = 6$};
  \node[rectangle, rounded corners=2pt, draw=violet!80!black, fill=violet!15, font=\scriptsize\bfseries] (tr) at (4.8,0) {3. HÀNG TRĂM\\$3 + 4 = 7$};

  \draw[arr, thick, red!85!black] (0.9,0) -- (1.5,0);
  \draw[arr, thick, red!85!black] (3.3,0) -- (3.9,0);
  \node[font=\tiny\bfseries, color=red!80!black] at (2.4,-0.6) {Tính từ PHẢI sang TRÁI};
\end{tikzpicture}
\end{tcolorbox}

\vspace{-1.5mm}
\noindent
\begin{minipage}[t]{0.49\textwidth}\vspace{0pt}
\begin{tcolorbox}[
  enhanced,
  colback=green!4!white,
  colframe=green!60!black,
  arc=2mm, boxrule=0.8pt,
  left=2mm, right=2mm, top=1.8mm, bottom=1.8mm,
  title={\footnotesize\bfseries\color{white}\faCheckCircle\ 4. Các dạng toán rèn luyện}
]
{\footnotesize
\textbf{1. Cộng số có 3 chữ số với số có 2 hoặc 1 chữ số:}\\
- $425 + 43 = \mathbf{468}$ (chú ý đặt $43$ thẳng cột chục và đơn vị).\\
- $651 + 7 = \mathbf{658}$ (đặt $7$ thẳng cột đơn vị).\\
\textbf{2. Tính nhẩm nhanh số tròn trăm, tròn chục:}\\
- $300 + 400 = \mathbf{700}$.\\
- $520 + 60 = \mathbf{580}$.\\
- $240 + 300 = \mathbf{540}$.
}
\end{tcolorbox}
\end{minipage}
\hfill
\begin{minipage}[t]{0.49\textwidth}\vspace{0pt}
\begin{tcolorbox}[
  enhanced,
  colback=rose!4!white,
  colframe=rose!70!black,
  arc=2mm, boxrule=0.8pt,
  left=2mm, right=2mm, top=1.8mm, bottom=1.8mm,
  title={\footnotesize\bfseries\color{white}\faExclamationTriangle\ 5. Lỗi đặt tính lệch cột nghiêm trọng}
]
{\footnotesize
\begin{itemize}\setlength{\itemsep}{1.5pt}\setlength{\leftskip}{-2.5mm}
  \item \textbf{Cực kì nguy hiểm:} Khi cộng $425 + 43$, viết chữ số $4$ của $43$ thẳng dưới số $4$ hàng trăm của $425$ là \textbf{sai hoàn toàn}!
  \item Chữ số $4$ của $43$ là hàng chục, phải thẳng với số $2$ hàng chục của $425$.
  \item Dùng thước kẻ kẻ vạch ngang thẳng tắp, không dùng tay kẻ nguệch ngoạc.
\end{itemize}
}
\end{tcolorbox}
\end{minipage}


\end{document}
